All experiments use the ISIC-2019 training collection, a public set of 25{,}331
dermoscopic images labelled across the eight diagnostic classes listed in
Chapter~1. The collection is itself a union of three sources HAM10000,
BCN20000, and MSK which we infer from the image-ID prefix, because keeping
track of the source matters for building an honest train/test split. Each image
also carries clinical metadata: approximate age, sex, the general anatomical
site of the lesion, and a lesion identifier that groups multiple images of the
same physical lesion.

The defining property of this dataset, and the one that drives most of our
design decisions, is its class imbalance. Table~\ref{tab:classdist} shows the
per-class counts for the full collection. The largest class, nevus, is more than
fifty times larger than the smallest, dermatofibroma. Several of the rare
classes (for example SCC and AK) are clinically important, so a model that
simply learns the head of this distribution would be useless in practice. This
is the central reason we report balanced multi-class accuracy rather than overall
accuracy throughout.

\begin{table}[!t]
\centering
\caption{Class distribution of the full ISIC-2019 collection (before
deduplication) and imbalance relative to the largest class (NV).}
\label{tab:classdist}
\begin{tabular}{llrr}
\toprule
Code & Full name & Images & Ratio vs.\ NV \\
\midrule
MEL  & Melanoma                  & 4{,}522  & 2.8$\times$  \\
NV   & Melanocytic nevus         & 12{,}875 & 1$\times$    \\
BCC  & Basal cell carcinoma      & 3{,}323  & 3.9$\times$  \\
AK   & Actinic keratosis         & 867      & 14.8$\times$ \\
BKL  & Benign keratosis          & 2{,}624  & 4.9$\times$  \\
DF   & Dermatofibroma            & 239      & 53.9$\times$ \\
VASC & Vascular lesion           & 253      & 50.9$\times$ \\
SCC  & Squamous cell carcinoma   & 628      & 20.5$\times$ \\
\bottomrule
\end{tabular}
\end{table}

The metadata is genuinely incomplete: roughly 30\% of age values, about 20\% of
sex values, and about 15\% of anatomical-site values are missing, with the
recorded ages centred near 54 years. Rather than fill these gaps with guessed
values, our requirement is that the model treat ``missing'' as a first-class,
learnable signal, which is reflected directly in the architecture (Chapter~4).
